\ifdefined\gkver\else\def\gkver{student}\fi
\documentclass[\gkver]{gaokaozhenti}
\newcommand{\ccnum}[1]{{\fontspec{Noto Serif CJK SC}#1}}
\begin{document}

\chapter{1992年上海}

\begin{enumerate}
\item
%% number: 1
%% paperName: 1992年上海高考第 1 题 4 分
%% typeId: 1
%% type: 单选题
%% chapter: 原子和原子核
%% point: 天然放射性
%% method: 无
%% score: 4
%% degree: 850
%% duplicateId: 0
%% body:
\ce{^{238}_{92}U} 经 3 次 $\alpha$ 衰变和 2 次 $\beta$ 衰变，最后得到的原子核中的质子数为\xzanswer{B}
\fourchoices[answer=B]
{92}
{88}
{138}
{226}

%% answer: B
%% memo:
\memoanswer{本题原卷及材料中未提供详解，待补充。}

\item
%% number: 2
%% paperName: 1992年上海高考第 2 题 4 分
%% typeId: 1
%% type: 单选题
%% chapter: 电路
%% point: 串并联电路
%% method: 无
%% score: 4
%% degree: 650
%% duplicateId: 0
%% body:
电阻 $R_{1}$ 和 $R_{2}$ 并联在电路中时，通过 $R_{1}$ 的电流强度是通过 $R_{2}$ 的电流强度的 $n$ 倍，则当 $R_{1}$ 和 $R_{2}$ 串联在电路中时，$R_{1}$ 两端的电压 $U_{1}$ 与 $R_{2}$ 两端的电压 $U_{2}$ 之比 $U_{1}/U_{2}$ 为\xzanswer{C}
\fourchoices[answer=C]
{$n$}
{$n^{2}$}
{$\frac{1}{n}$}
{$\frac{1}{n^{2}}$}

%% answer: C
%% memo:
\memoanswer{本题原卷及材料中未提供详解，待补充。}

\item
%% number: 3
%% paperName: 1992年上海高考第 3 题 4 分
%% typeId: 1
%% type: 单选题
%% chapter: 交流电
%% point: 变压器
%% method: 无
%% score: 4
%% degree: 650
%% duplicateId: 0
%% body:
下列四图中，a、b 为输入端，接交流电源，c、d 为输出端，则输出电压大于输入电压的电路是\xzanswer{C}
\fourchoices[answer=C, ispicture=true, h=2.2cm]
{\includesvg{figs/03a}}
{\includesvg{figs/03b}}
{\includesvg{figs/03c}}
{\includesvg{figs/03d}}

%% answer: C
%% memo:
\memoanswer{本题原卷及材料中未提供详解，待补充。}

\item
%% number: 4
%% paperName: 1992年上海高考第 4 题 4 分
%% typeId: 1
%% type: 单选题
%% chapter: 原子和原子核
%% point: 玻尔模型
%% method: 无
%% score: 4
%% degree: 650
%% duplicateId: 0
%% body:
根据玻尔理论，在氢原子中，量子数 $n$ 越大，则\xzanswer{B}
\fourchoices[answer=B]
{电子轨道半径越小}
{核外电子速度越小}
{原子能级的能量越小}
{原子的电势能越小}

%% answer: B
%% memo:
\memoanswer{本题原卷及材料中未提供详解，待补充。}

\item
%% number: 5
%% paperName: 1992年上海高考第 5 题 4 分
%% typeId: 1
%% type: 单选题
%% chapter: 电磁感应
%% point: 电磁感应能量分析
%% method: 无
%% score: 4
%% degree: 650
%% duplicateId: 0
%% body:
如图所示，闭合导线框的质量可以忽略不计，将它从图示位置匀速拉出匀强磁场。若第一次用 $0.3\Us$ 时间拉出，外力所做的功为 $W_{1}$，通过导线截面的电量为 $q_{1}$；第二次用 $0.9\Us$ 时间拉出，外力所做的功为 $W_{2}$，通过导线截面的电量为 $q_{2}$，则\xzanswer{C}
\onepicture[w=4.5cm]{\includesvg{figs/05}}
\fourchoices[answer=C]
{$W_{1} < W_{2}$，$q_{1} < q_{2}$}
{$W_{1} < W_{2}$，$q_{1} = q_{2}$}
{$W_{1} > W_{2}$，$q_{1} = q_{2}$}
{$W_{1} > W_{2}$，$q_{1} > q_{2}$}

%% answer: C
%% memo:
\memoanswer{本题原卷及材料中未提供详解，待补充。}

\item
%% number: 6
%% paperName: 1992年上海高考第 6 题 4 分
%% typeId: 1
%% type: 单选题
%% chapter: 机械波
%% point: 波长频率波速
%% method: 无
%% score: 4
%% degree: 450
%% duplicateId: 0
%% body:
如图所示，S 点为振源，其频率为 $100\UHz$，所产生的横波向右传播，波速为 $80\Ums$。P、Q 是波传播途径中的两点，已知 SP $=$ $4.2\Um$，SQ $=$ $5.4\Um$。当 S 通过平衡位置向上运动时，则 PQ\xzanswer{A}
\onepicture[w=8cm]{\includesvg{figs/06}}
\fourchoices[answer=A]
{P 在波谷，Q 在波峰}
{P 在波峰，Q 在波谷}
{P、Q 都在波峰}
{P 通过平衡位置向上运动，Q 通过平衡位置向下运动}

%% answer: A
%% memo:
\memoanswer{本题原卷及材料中未提供详解，待补充。}

\item
%% number: 7
%% paperName: 1992年上海高考第 7 题 4 分
%% typeId: 1
%% type: 单选题
%% chapter: 运动和力
%% point: 超重失重
%% method: 整体与隔离
%% score: 4
%% degree: 650
%% duplicateId: 0
%% body:
图中 A 为电磁铁，C 为胶木秤盘，A 和 C（包括支架）的总质量为 $M$，B 为铁片，质量为 $m$，整个装置用轻绳悬挂于 O 点。当电磁铁通电，铁片被吸引上升的过程中，轻绳上拉力 $F$ 的大小为\xzanswer{D}
\onepicture[h=4cm]{\includesvg{figs/07}}
\fourchoices[answer=D]
{$F = Mg$}
{$Mg < F < (M + m)g$}
{$F = (M + m)g$}
{$F > (M + m)g$}

%% answer: D
%% memo:
\memoanswer{本题原卷及材料中未提供详解，待补充。}

\item
%% number: 8
%% paperName: 1992年上海高考第 8 题 4 分
%% typeId: 1
%% type: 单选题
%% chapter: 光学
%% point: 光的折射
%% method: 无
%% score: 4
%% degree: 650
%% duplicateId: 0
%% body:
右图表示一条光线由空气射到半圆形玻璃砖表面的圆心 O 处，在玻璃砖的半圆形表面镀有银反射面，则以下几个光路图中，正确、完整表示光线行进过程的是\xzanswer{D}
\onepicture[w=3.2cm]{\includesvg{figs/08}}
\fourchoices[answer=D, ispicture=true, h=2.6cm]
{\includesvg{figs/08a}}
{\includesvg{figs/08b}}
{\includesvg{figs/08c}}
{\includesvg{figs/08d}}

%% answer: D
%% memo:
\memoanswer{本题原卷及材料中未提供详解，待补充。}

\item
%% number: 9
%% paperName: 1992年上海高考第 9 题 5 分
%% typeId: 2
%% type: 多选题
%% chapter: 电磁感应
%% point: 楞次定律推广
%% method: 无
%% score: 5
%% degree: 650
%% duplicateId: 0
%% body:
在水平面上有一固定的 U 形金属框架，框架上置一金属杆 ab，如图所示（纸面即水平面）。在垂直纸面方向有一匀强磁场，则\xzanswer{BD}
\onepicture[w=3.2cm]{\includesvg{figs/09}}
\fourchoices[answer=BD]
{若磁场方向垂直纸面向外并增长时，杆 ab 将向右移动}
{若磁场方向垂直纸面向外并减少时，杆 ab 将向右移动}
{若磁场方向垂直纸面向里并增长时，杆 ab 将向右移动}
{若磁场方向垂直纸面向里并减少时，杆 ab 将向右移动}

%% answer: BD
%% memo:
\memoanswer{本题原卷及材料中未提供详解，待补充。}

\item
%% number: 10
%% paperName: 1992年上海高考第 10 题 5 分
%% typeId: 1
%% type: 单选题
%% chapter: 固体液体的性质
%% point: 表面张力
%% method: 无
%% score: 5
%% degree: 650
%% duplicateId: 0
%% body:
下列说法正确的是\xzanswer{D}
\fourchoices[answer=D]
{质量相同、温度相同的两个物体的内能一定相同}
{一定量气体的体积由 $10\UL$ 膨胀为 $20\UL$ 时，其压强一定变为原来的一半}
{大多数金属都是各向同性的，它们都是非晶体}
{液体表面具有收缩的趋势，是由于在液体表面层里分子的分布比内部稀疏的缘故}

%% answer: D
%% memo:
\memoanswer{本题原卷及材料中未提供详解，待补充。}

\item
%% number: 11
%% paperName: 1992年上海高考第 11 题 5 分
%% typeId: 2
%% type: 多选题
%% chapter: 功和能
%% point: 动能定理
%% method: 无
%% score: 5
%% degree: 650
%% duplicateId: 0
%% body:
在粗糙水平面上运动的物体，从 A 点开始受水平恒力 $F$ 作用作直线运动到 B 点。已知物体在 B 点的速度与在 A 点的速度大小相等，则在这过程中\xzanswer{AC}
\fourchoices[answer=AC]
{物体不一定作匀速直线运动}
{$F$ 始终与摩擦力方向相反}
{$F$ 与摩擦力对物体所作总功为零}
{$F$ 与摩擦力对物体的总冲量为零}

%% answer: AC
%% memo:
\memoanswer{本题原卷及材料中未提供详解，待补充。}

\item
%% number: 12
%% paperName: 1992年上海高考第 12 题 5 分
%% typeId: 2
%% type: 多选题
%% chapter: 交流电
%% point: 交流电
%% method: 无
%% score: 5
%% degree: 650
%% duplicateId: 0
%% body:
将阻值为 $R$ 的电阻，接在电动势 $\varepsilon = 10\UV$，内电阻 $r = R$ 的直流电源两端，此时电阻 $R$ 上的电功率为 $P$\xzanswer{BC}
\fourchoices[answer=BC]
{若将电阻接到直流电源上，电阻两端电压 $U = 10\UV$，则其电功率为 $2P$}
{若将电阻接到直流电源上，电阻两端电压 $U = 20\UV$，则其电功率为 $16P$}
{若将电阻接到正弦交流电源上，电阻两端电压的最大值 $U_{m} = 10\UV$，则其电功率为 $2P$}
{若将电阻接到正弦交流电源上，电阻两端电压的最大值 $U_{m} = 20\UV$，则其电功率为 $4P$}

%% answer: BC
%% memo:
\memoanswer{本题原卷及材料中未提供详解，待补充。}

\item
%% number: 13
%% paperName: 1992年上海高考第 13 题 5 分
%% typeId: 2
%% type: 多选题
%% chapter: 功和能
%% point: 功
%% method: 无
%% score: 5
%% degree: 450
%% duplicateId: 0
%% body:
质量为 $m$ 的物块始终固定在倾角为 $\theta$ 的斜面上，下列说法中正确的是\xzanswer{ABC}
\onepicture[w=4.5cm]{\includesvg{figs/13}}
\fourchoices[answer=ABC]
{若斜面向右匀速移动距离 $s$，斜面对物块没有做功}
{若斜面向上匀速移动距离 $s$，斜面对物块做功 $mgs$}
{若斜面向左以加速度 $a$ 移动距离 $s$，斜面对物块做功 $mas$}
{若斜面向下以加速度 $a$ 移动距离 $s$，斜面对物块做功 $m(g + a)s$}

%% answer: ABC
%% memo:
\memoanswer{本题原卷及材料中未提供详解，待补充。}

\item
%% number: 14
%% paperName: 1992年上海高考第 14 题 4 分
%% typeId: 3
%% type: 填空题
%% chapter: 光学
%% point: 光电效应
%% method: 无
%% score: 4
%% degree: 650
%% duplicateId: 0
%% body:
图示为光电管的工作电路，则图中电源的正极为\tkanswer[1.5cm]{a}（填 a 或 b）。若使这种光电管产生光电效应的入射光的最大波长为 $\lambda$，则能使光电管工作的入射光光子的最小能量为\tkanswer[3cm]{$hc/\lambda$}。
\onepicture[align=right, h=4cm]{\includesvg{figs/14}}

%% answer: a，$hc/\lambda$
\jdanswer{
a，$hc/\lambda$
}

%% memo:
\memoanswer{本题原卷及材料中未提供详解，待补充。}

\item
%% number: 15
%% paperName: 1992年上海高考第 15 题 4 分
%% typeId: 3
%% type: 填空题
%% chapter: 气体性质
%% point: 用单分子油膜估测分子大小
%% method: 无
%% score: 4
%% degree: 650
%% duplicateId: 0
%% body:
将 $1\Ucmc$ 的油酸溶于酒精，制成 $200\Ucmc$ 的油酸酒精溶液。已知 $1\Ucmc$ 溶液有 50 滴，现取 1 滴油酸酒精溶液滴到水面上。随着酒精溶于水，油酸在水面上形成一单分子薄层，已测出这一薄层的面积为 $0.2\Umq$。由此可估测油酸分子的直径为\tkanswer[2.6cm]{$5\times10^{-10}\Um$}。

%% answer: $5\times10^{-10}\Um$
\jdanswer{
$5\times10^{-10}\Um$
}

%% memo:
\memoanswer{本题原卷及材料中未提供详解，待补充。}

\item
%% number: 16
%% paperName: 1992年上海高考第 16 题 4 分
%% typeId: 3
%% type: 填空题
%% chapter: 功和能
%% point: 机械能守恒定律
%% method: 无
%% score: 4
%% degree: 650
%% duplicateId: 0
%% body:
如图所示，质量均为 $m$ 的小球 A、B、C，用两条长为 $l$ 的细线相连，置于高为 $h$ 的光滑水平桌面上，$l > h$，A 球刚跨过桌边。若 A 球、B 球相继下落着地后均不再反跳，则 C 球离开桌边时的速度大小是\tkanswer[2.6cm]{$\sqrt{\frac{5gh}{3}}$}。
\onepicture[align=right, w=6cm]{\includesvg{figs/16}}

%% answer: $\sqrt{\frac{5gh}{3}}$
\jdanswer{
$\sqrt{\frac{5gh}{3}}$
}

%% memo:
\memoanswer{本题原卷及材料中未提供详解，待补充。}

\item
%% number: 17
%% paperName: 1992年上海高考第 17 题 4 分
%% typeId: 3
%% type: 填空题
%% chapter: 电场
%% point: 带电粒子的圆周运动
%% method: 极值
%% score: 4
%% degree: 450
%% duplicateId: 0
%% body:
半径为 $r$ 的绝缘光滑圆环固定在竖直平面内，环上套有一质量为 $m$、带正电的珠子，空间存在水平向右的匀强电场，如图所示。珠子所受静电力是其重力的 $\frac{3}{4}$ 倍。将珠子从环上最低位置 A 点静止释放，则珠子所能获得的最大动能 $E_{k} =$\tkanswer[2cm]{$\frac{1}{4}mgr$}。
\onepicture[align=right, w=5cm]{\includesvg{figs/17}}

%% answer: $\frac{1}{4}mgr$
\jdanswer{
$\frac{1}{4}mgr$
}

%% memo:
\memoanswer{本题原卷及材料中未提供详解，待补充。}

\item
%% number: 18
%% paperName: 1992年上海高考第 18 题 4 分
%% typeId: 3
%% type: 填空题
%% chapter: 电路
%% point: 闭合电路欧姆定律
%% method: 无
%% score: 4
%% degree: 650
%% duplicateId: 0
%% body:
用量程为 $500\UuA$ 的电流表改装而成的、量程为 $10\UV$ 的电压表，与 $10\UkO$ 的电阻串联后，接在一电源的两端，这时伏特表的读数为 $8\UV$。忽略电源内阻，则电源的电动势 $\varepsilon =$\tkanswer[1.8cm]{$12\UV$}。

%% answer: $12\UV$
\jdanswer{
$12\UV$
}

%% memo:
\memoanswer{本题原卷及材料中未提供详解，待补充。}

\item
%% number: 19
%% paperName: 1992年上海高考第 19 题 4 分
%% typeId: 3
%% type: 填空题
%% chapter: 电磁感应
%% point: 电磁感应能量分析
%% method: 无
%% score: 4
%% degree: 650
%% duplicateId: 0
%% body:
如图所示，MN 为金属杆，在竖直平面内贴着光滑金属导轨下滑，导轨的间距 $l = 10\Ucm$，导轨上端接有电阻 $R = 0.5\UO$，导轨与金属杆电阻不计，整个装置处于 $B = 0.5\UT$ 的水平匀强磁场中。若杆稳定下落时，每秒钟有 $0.02\UJ$ 的重力势能转化为电能，则 MN 杆的下落速度 $v =$ \tkanswer[1.8cm]{$2\Ums$}。
\onepicture[align=right, h=4.5cm]{\includesvg{figs/19}}

%% answer: $2\Ums$
\jdanswer{
$2\Ums$
}

%% memo:
\memoanswer{本题原卷及材料中未提供详解，待补充。}

\item
%% number: 20
%% paperName: 1992年上海高考第 20 题 4 分
%% typeId: 3
%% type: 填空题
%% chapter: 运动和力
%% point: 动量定理
%% method: 无
%% score: 4
%% degree: 650
%% duplicateId: 0
%% body:
两物体的质量为 $m_{1}$ 和 $m_{2}$，它们分别在恒力 $F_{1}$ 和 $F_{2}$ 的作用下由静止开始运动，经相同的位移，动量的增加量相同，则两恒力的比值 $F_{1}/F_{2} =$\tkanswer[2cm]{$m_{2}/m_{1}$}。

%% answer: $m_{2}/m_{1}$
\jdanswer{
$m_{2}/m_{1}$
}

%% memo:
\memoanswer{本题原卷及材料中未提供详解，待补充。}

\item
%% number: 21
%% paperName: 1992年上海高考第 21 题 4 分
%% typeId: 3
%% type: 填空题
%% chapter: 光学
%% point: 透镜
%% method: 无
%% score: 4
%% degree: 450
%% duplicateId: 0
%% body:
某人透过焦距为 $6\Ucm$，直径为 $3\Ucm$ 的凸透镜看报，将离眼 $16\Ucm$ 处的报纸成象在离眼 $24\Ucm$ 处。为此，应使透镜与报纸相距\tkanswer[1.5cm]{$4\Ucm$}。设眼在透镜主轴上，报纸平面垂直于主轴，若报上密排着宽、高均为 $0.3\Ucm$ 的字，则他通过透镜至多能看清同一行上\tkanswer[1.2cm]{6}个完整的字（忽略眼睛瞳孔的大小）。

%% answer: $4\Ucm$，6
\jdanswer{
$4\Ucm$，6
}

%% memo:
\memoanswer{本题原卷及材料中未提供详解，待补充。}

\item
%% number: 22
%% paperName: 1992年上海高考第 22 题 4 分
%% typeId: 4
%% type: 实验题
%% chapter: 气体性质
%% point: 分子动理论
%% method: 无
%% score: 4
%% degree: 650
%% duplicateId: 0
%% body:
关于布朗运动的实验，下列说法正确的是\xzanswer{D}
\onepicture[w=5.5cm]{\includesvg{figs/22}}
\fourchoices[answer=D]
{图中记录的是分子无规则运动的情况}
{图中记录的是微粒作布朗运动的轨迹}
{实验中可以看到，微粒越大，布朗运动越明显}
{实验中可以看到，温度越高，布朗运动越激烈}

%% answer: D
\jdanswer{
D
}
%% memo:
\memoanswer{本题原卷及材料中未提供详解，待补充。}

\item
%% number: 23
%% paperName: 1992年上海高考第 23 题 5 分
%% typeId: 4
%% type: 实验题
%% chapter: 电路
%% point: 多用表的使用
%% method: 无
%% score: 5
%% degree: 650
%% duplicateId: 0
%% body:
在下列实验中，哪些操作是错误的\xzanswer{CD}
\fourchoices[answer=CD]
{在《验证机械能守恒定律》实验中，应从几条打点纸带中挑出第一、二两点间的距离近似等于 $2\Umm$，且点迹清晰的纸带来测算}
{在用公式 $f = \frac{L^{2} - d^{2}}{4L}$ 求焦距的实验中，应先固定光源与光屏的位置，并使它们之间的距离 $L > 4f$}
{在使用欧姆表时，先要进行欧姆档调零。在以后操作中，不论被测电阻的阻值是多少，均不必再次进行欧姆档调零}
{在《用单摆测定重力加速度》的实验中，使摆球在摆角小于 $5^{\circ}$ 时从静止释放，并同时开始计时}

%% answer: CD
\jdanswer{
CD
}
%% memo:
\memoanswer{本题原卷及材料中未提供详解，待补充。}

\item
%% number: 24
%% paperName: 1992年上海高考第 24 题 5 分
%% typeId: 4
%% type: 实验题
%% chapter: 电磁感应
%% point: 自感与互感，涡流
%% method: 无
%% score: 5
%% degree: 650
%% duplicateId: 0
%% body:
在图示实验中，带铁芯的、电阻较小的线圈 L 与灯 A 并联。当合上电键 K，灯 A 正常发光。试判断下列说法中哪些是正确的\xzanswer{BC}
\onepicture[w=4.5cm]{\includesvg{figs/24}}
\fourchoices[answer=BC]
{当断开 K 时，灯 A 立即熄灭}
{当断开 K 时，灯 A 突然闪亮后熄灭}
{若用阻值与线圈 L 相同的电阻取代 L 接入电路，当断开 K，灯 A 立即熄灭}
{若用阻值与线圈 L 相同的电阻取代 L 接入电路，当断开 K，灯 A 突然闪亮后熄灭}

%% answer: BC
\jdanswer{
BC
}
%% memo:
\memoanswer{本题原卷及材料中未提供详解，待补充。}

\item
%% number: 25
%% paperName: 1992年上海高考第 25 题 4 分
%% typeId: 4
%% type: 实验题
%% chapter: 光学
%% point: 测定玻璃折射率
%% method: 无
%% score: 4
%% degree: 650
%% duplicateId: 0
%% body:
在《测定玻璃的折射率》的实验中，已画好玻璃砖界面两直线 $aa'$ 和 $bb'$ 后，不觉误将玻璃砖向上平移如图中虚线所示，若其他操作正确，则测得的折射率将\tkanswer[1.5cm]{不变}（填“偏大”、“偏小”或“不变”）。
\onepicture[align=right, w=5.5cm]{\includesvg{figs/25}}

%% answer: 不变
\jdanswer{
不变
}

%% memo:
\memoanswer{本题原卷及材料中未提供详解，待补充。}

\item
%% number: 26
%% paperName: 1992年上海高考第 26 题 6 分
%% typeId: 4
%% type: 实验题
%% chapter: 运动和力
%% point: 动量守恒定律
%% method: 无
%% score: 6
%% degree: 450
%% duplicateId: 0
%% body:
某同学设计了一个用打点计时器验证动量守恒定律的实验：在小车 A 的前端粘有橡皮泥，推动小车 A 使之作匀速运动，然后与原来静止在前方的小车 B 相碰并粘合成一体，继续作匀速运动。他设计的具体装置如图所示，在小车 A 后连着纸带，电磁打点计时器电源频率为 $50\UHz$，长木板下垫着小木片用以平衡摩擦力。
\begin{enumerate}
	\item
	若已得到打点纸带如上图，并测得各计数点间距标在图上。A 为运动起始的第一点。则应选\tkanswer[1.2cm]{BC}段来计算 A 的碰前速度。应选\tkanswer[1.2cm]{DE}段来计算 A 和 B 碰后的共同速度。（以上两格填“AB”或“BC”或“CD”或“DE”）。
	\item
	已测得小车 A 的质量 $m_{1} = 0.40\Ukg$，小车 B 的质量 $m_{2} = 0.20\Ukg$。由以上测量结果可得：碰前总动量 $=$\tkanswer[2.2cm]{$0.42\Ukgms$}；碰后总动量 $=$\tkanswer[2.2cm]{$0.417\Ukgms$}（或 $0.42\Ukgms$）。
\end{enumerate}
\onepicture[align=right, w=12cm]{\includesvg{figs/26}}

%% answer: （1）BC，DE；（2）$0.42\Ukgms$，$0.417\Ukgms$（或 $0.42\Ukgms$）
\jdanswer{
\begin{enumerate}
	\item
	BC，DE

	\item
	$0.42\Ukgms$，$0.417\Ukgms$（或 $0.42\Ukgms$）

\end{enumerate}
}

%% memo:
\memoanswer{本题原卷及材料中未提供详解，待补充。}

\item
%% number: 27
%% paperName: 1992年上海高考第 27 题 10 分
%% typeId: 6
%% type: 计算题
%% chapter: 气体性质
%% point: 玻意耳定律
%% method: 无
%% score: 10
%% degree: 450
%% duplicateId: 0
%% body:
某水银气压计的玻璃管顶端高出水银槽液面 $1\Um$，因上部混入少量空气，使其读数不准，当气温为 $27\UdC$，标准气压计读数为 $76\UcmHg$ 时，该气压计读数为 $70\UcmHg$。
\begin{enumerate}
	\item
	在相同气温下，若用该气压计测量气压，测得读数为 $68\UcmHg$，则实际气压应为多少 $\UcmHg$？
	\item
	若在气温为 $-3\UdC$ 时，用该气压计测得气压读数仍为 $70\UcmHg$，则实际气压应为多少 $\UcmHg$？
\end{enumerate}
\onepicture[align=right, h=5cm]{\includesvg{figs/27}}

%% answer: （1）$73.6\UcmHg$；（2）$75.4\UcmHg$
\jdanswer{
\begin{enumerate}
	\item
	$73.6\UcmHg$

	\item
	$75.4\UcmHg$

\end{enumerate}
}

%% memo:
\memoanswer{
（1）以气压计内混入气体为考察对象，因温度不变，有

$p_{0}V_{0} = p_{1}V_{1}$　　　　①

设管截面积为 $S$，压强以 $\UcmHg$ 为单位，体积以 $\Ucmc$ 为单位，则

$p_{0} = 76 - 70 = 6$，$V_{0} = (100 - 70)S = 30S$，$h = (100 - 68)S = 32S$

代入①式，得 $p_{1} = \frac{30}{32}\times6\UcmHg = 5.625\UcmHg$

实际气压

$p = p_{1} + 68 = 73.625\UcmHg$

（2）因体积不变，有

$\frac{p_{0}}{T_{0}} = \frac{p_{2}}{T_{2}}$　　　　②

以 $p_{0} = 6$，$T_{0} = 273 + 27 = 300\UK$，$T_{2} = 273 - 3 = 270\UK$

代入②式，得：$p_{2} = \frac{270}{300}\times6\UcmHg = 5.4\UcmHg$

实际气压：$p' = p_{2} + 70 = 75.4\UcmHg$

评分标准：本题共 10 分。第（1）小题 5 分，第（2）小题 5 分。

（1）写出①式，并正确表示各量得 3 分。正确得出 $p_{1}$ 结果得 1 分，正确得出实际气压得 1 分。

（2）写出②式，并正确表示各量得 3 分。正确得出 $p_{2}$ 结果得 1 分，正确得出实际气压得 1 分。
}

\item
%% number: 28
%% paperName: 1992年上海高考第 28 题 12 分
%% typeId: 6
%% type: 计算题
%% chapter: 运动和力
%% point: 动量守恒定律
%% method: 无
%% score: 12
%% degree: 450
%% duplicateId: 0
%% body:
\begin{enumerate}
	\item
	在光滑水平台面上，质量 $m_{1} = 4\Ukg$ 的物块 1 具有动能 $E = 100\UJ$。物块 1 与原来静止的质量 $m_{2} = 1\Ukg$ 的物块 2 发生碰撞，碰后粘合在一起。求碰撞中的机械能损失 $\Delta E$。
	\item
	若物块 1、2 分别具有动能 $E_{1}$、$E_{2}$，$E_{1}$ 与 $E_{2}$ 之和为 $100\UJ$。两物块相向运动而碰撞并粘合在一起，问 $E_{1}$、$E_{2}$ 各应为多少时，碰撞中损失的机械能最大？（需说明理由。）这时损失的机械能为多少 $\Delta E'$？
\end{enumerate}

%% answer: （1）$20\UJ$；（2）$100\UJ$
\jdanswer{
\begin{enumerate}
	\item
	$20\UJ$

	\item
	$100\UJ$

\end{enumerate}
}

%% memo:
\memoanswer{
解法一

（1）物块 1 碰前速度：$v_{1} = \sqrt{\frac{2E}{m_{1}}} = \sqrt{50}\Ums$　　　　①

由动量守恒：$m_{1}v_{1} = (m_{1} + m_{2})v'$　　　　②

碰后共同速度：$v' = \frac{m_{1}v_{1}}{m_{1} + m_{2}}$　　　　③

碰后总动能：$E' = \frac{1}{2}(m_{1} + m_{2}){v'}^{2} = \frac{m_{1}^{2}v_{1}^{2}}{2(m_{1} + m_{2})}$　　　　④

代入数值　　$E' = 80\UJ$

$\Delta E = E - E' = 20\UJ$　　　　⑤

（2）设碰前物块 1、2 速度为 $v_{1}$、$v_{2}$，碰后共同速度为 $v''$，由动量守恒，有

$m_{1}v_{1}' - m_{2}v_{2}' = (m_{1} + m_{2})v''$　　　　⑥

碰后两物块总动能　　$E'' = \frac{1}{2}(m_{1} + m_{2}){v''}^{2}$

以⑥式代入　　$E'' = \frac{1}{2}\frac{(m_{1}v_{1}' - m_{2}v_{2}')^{2}}{m_{1} + m_{2}}$　　　　⑦

由⑦式，当　　$m_{1}v_{1}' = m_{2}v_{2}'$　　　　⑧

时最小，其值为零，其时机械能损失最大。以 $m_{1}$、$m_{2}$ 值代入⑧式，有

$4v_{1}' = v_{2}'$　　　　⑨

由题意　　$\frac{1}{2}m_{1}v_{1}'^{2} + \frac{1}{2}m_{2}v_{2}'^{2} = 100$

以⑨式代入，解得

$v_{1}' = \sqrt{10}\Ums$，$v_{2}' = 4\sqrt{10}\Ums$

$E_{1} = \frac{1}{2}m_{1}v_{1}'^{2} = 20\UJ$　　　　⑩

$E_{2} = \frac{1}{2}m_{2}v_{2}'^{2} = 80\UJ$　　　　\ccnum{⑪}

$\Delta E' = E - E'' = 100\UJ$　　　　\ccnum{⑫}

解法二

（1）物块 1 动量　　$p = \sqrt{2m_{1}E_{k}}$　　　　①

碰后动量守恒，故碰后两物块总动能

$E' = \frac{p^{2}}{2(m_{1} + m_{2})} = \frac{2m_{1}E}{2(m_{1} + m_{2})} = \frac{2\times4\times100}{2(4 + 1)}\UJ = 80\UJ$　　　　②

$\Delta E = E - E' = 20\UJ$　　　　③

（2）碰前物块 1、2 动量　　$p_{1} = \sqrt{2m_{1}E_{1}}$，$p_{2} = \sqrt{2m_{2}E_{2}}$　　　　④

两物块总动量　　$p = \sqrt{2m_{1}E_{1}} - \sqrt{2m_{2}E_{2}}$　　　　⑤

由动量守恒，碰后两物块总动能

$E'' = \frac{(p_{1} - p_{2})^{2}}{2(m_{1} + m_{2})} = \frac{(\sqrt{2m_{1}E_{1}} - \sqrt{2m_{2}E_{2}})^{2}}{2(m_{1} + m_{2})}$　　　　⑥

由⑥式，当

$p_{1} = p_{2}$ 或 $\sqrt{2m_{1}E_{1}} = \sqrt{2m_{2}E_{2}}$

即　　$m_{1}E_{1} = m_{2}E_{2}$　　　　⑦

时 $E''$ 最小，其值为 0，其时机械能损失最大。

由⑦式及　　$E_{1} + E_{2} = 100\UJ$

即可解得　　$E_{1} = 20\UJ$，$E_{2} = 80\UJ$

$\Delta E' = E - E'' = 100\UJ$

评分标准：本题共 12 分。第（1）小题 5 分，第（2）小题 7 分。

（1）利用动量守恒正确得出 $E'$ 结果得 4 分，正确得出 $\Delta E$ 结果得 1 分。

（2）正确得出 $\Delta E'$ 最大条件（$m_{1}v_{1}' = m_{2}v_{2}'$ 或 $m_{1}E_{1} = m_{2}E_{2}$）得 2 分，正确说明理由（其时 $E''$ 最小）得 2 分，正确得出 $E_{1}$、$E_{2}$ 结果各得 1 分，正确得出 $\Delta E'$ 结果得 1 分。
}

\item
%% number: 29
%% paperName: 1992年上海高考第 29 题 15 分
%% typeId: 6
%% type: 计算题
%% chapter: 磁场
%% point: 带电粒子在磁场中的运动
%% method: 无
%% score: 15
%% degree: 450
%% duplicateId: 0
%% body:
图示为一种获得高能粒子的装置。环形区域内存在垂直纸面向外、大小可调节的均匀磁场。质量为 $m$，电量为 $+q$ 的粒子在环中作半径为 $R$ 的圆周运动。A、B 为两块中心开有小孔的极板，原来电势都为零，每当粒子飞经 A 板时，A 板电势升高为 $+U$，B 板电势仍保持为零，粒子在两板间电场中得到加速。每当粒子离开 B 板时，A 板电势又降为零。粒子在电场一次次加速下动能不断增大，而绕行半径不变。
\begin{enumerate}
	\item
	设 $t = 0$ 时粒子静止在 A 板小孔处，在电场作用下加速，并绕行第一圈。求粒子绕行 $n$ 圈回到 A 板时获得的总动能 $E_{n}$。
	\item
	为使粒子始终保持在半径为 $R$ 的圆轨道上运动，磁场必须周期性递增。求粒子绕行第 $n$ 圈时的磁感应强度 $B_{n}$。
	\item
	求粒子绕行 $n$ 圈所需的总时间 $t_{n}$（设极板间距远小于 $R$）。
	\item
	在右下图中画出 A 板电势 $u$ 与时间 $t$ 的关系（从 $t = 0$ 起画到粒子第四次离开 B 板时即可）。
	\item
	在粒子绕行的整个过程中，A 板电势是否可始终保持为 $+U$？为什么？
\end{enumerate}
\twopicture[labelA=1992上海29a, labelB=1992上海29b, align=right, wA=4.2cm, wB=4.8cm]
{\includesvg{figs/29a}}
{\includesvg{figs/29b}}

%% answer: （1）$E_{n} = nqU$；（2）$B_{n} = \frac{1}{R}\sqrt{\frac{2nmU}{q}}$；（3）$t_{n} = 2\pi R\sqrt{\frac{m}{2qU}}\left(1 + \frac{1}{\sqrt{2}} + \frac{1}{\sqrt{3}} + \cdots + \frac{1}{\sqrt{n}}\right)$；（4）如下图；（5）不可以
\jdanswer{
\begin{enumerate}
	\item
	$E_{n} = nqU$

	\item
	$B_{n} = \frac{1}{R}\sqrt{\frac{2nmU}{q}}$

	\item
	$t_{n} = 2\pi R\sqrt{\frac{m}{2qU}}\left(1 + \frac{1}{\sqrt{2}} + \frac{1}{\sqrt{3}} + \cdots + \frac{1}{\sqrt{n}}\right)$

	\item
	如下图。（对图的要求：间隔越来越近的等幅脉冲）

	\item
	不可以。因为这样粒子在 A、B 之间飞行时电场对其作功 $+qU$ 使之加速，在 A、B 之外飞行时电场又对其作功 $-qU$ 使之减速，粒子绕行一周电场对其所作总功为零，能量不会增大。

\end{enumerate}
}

%% memo:
\memoanswer{
（1）$E_{n} = nqU$

（2）$nqU = \frac{1}{2}mv_{n}^{2}$，　　$v_{n} = \sqrt{\frac{2nqU}{m}}$

$\frac{mv_{n}^{2}}{R} = qv_{n}B_{n}$，　　$B_{n} = \frac{mv_{n}}{qR}$

以 $v_{n}$ 的结果代入，　　$B_{n} = \frac{m}{qR}\sqrt{\frac{2nqU}{m}} = \frac{1}{R}\sqrt{\frac{2nmU}{q}}$

（3）绕行第 $n$ 圈需时　　$\frac{2\pi R}{v_{n}} = 2\pi R\sqrt{\frac{m}{2qU}}\frac{1}{\sqrt{n}}$

得　　$t_{n} = 2\pi R\sqrt{\frac{m}{2qU}}\left(1 + \frac{1}{\sqrt{2}} + \frac{1}{\sqrt{3}} + \cdots + \frac{1}{\sqrt{n}}\right)$

（4）如下图。（对图的要求：间隔越来越近的等幅脉冲）

\onepicture[w=5.5cm]{\includesvg{figs/29_an}}

（5）不可以。因为这样粒子在 A、B 之间飞行时电场对其作功 $+qU$ 使之加速，在 A、B 之外飞行时电场又对其作功 $-qU$ 使之减速，粒子绕行一周电场对其所作总功为零，能量不会增大。
}

\end{enumerate}

\end{document}
